% part: intuitionistic-logic
% chapter: introduction
% section: constructive-reasoning

\documentclass[../../../include/open-logic-chapter]{subfiles}

\begin{document}

\olfileid{int}{int}{cr}

\section{تعميري استدلال}

په مودالي عاملونو يا د دويمې درجې کميت‌ټاکونکو سره د کلاسيک منطق
د پراخولو پر خلاف، شهودي منطق په دې معنا «غيرکلاسيک» دے چې کلاسيک
منطق محدودوي۔ کلاسيک منطق له بېلابېلو اړخونو \emph{ناتعميري} دے۔
شهودي منطق غواړي هغه ډېر «تعميري» استدلال ونيسي چې د تعميري رياضياتو
له يوه ډول سره تړاو لري۔ لاندې بېلګې د دې نظر ځينې بنسټيزې انګېزې
روښانولے شي۔

فرض کړئ يو څوک دعوه وکړي چې داسې طبيعي عدد~$n$ ئې ټاکلے دے چې که
$n$ جوړ وي، د ريمان فرضيه رښتيا ده، او که $n$ طاق وي، د ريمان فرضيه
نارښتيا ده۔ دا خو زېرے شو!{} د ريمان د فرضيې رښتياوالی د رياضياتو له
سترو پرانيستو پوښتنو څخه دے، او ښکاري چې دغه کس ستونزه يوازې يوې
محاسبې ته راکمه کړې ده: دا معلومول چې يو ځانګړے عدد جوړ دے که نه۔

د~$n$ دا جادويي قيمت څو دے؟ هغه ئې داسې بيانوي: $n$ هغه طبيعي عدد
دے چې که د ريمان فرضيه رښتيا وي، له $2$ سره برابر دے، او که نه، له
$3$ سره برابر دے۔

تاسو په غصه خپلې پيسې بېرته غواړئ۔ له کلاسيک نظره پورته بيان په رښتيا
د $n$ يو يوازينے قيمت ټاکي؛ خو هغه څۀ چې تاسو ئې غواړئ، دا دي چې د
$n$ قيمت \emph{په څرګنده توګه} درکړے شي۔

يوه بله، ښايي لږه تصنعي بېلګه واخلئ۔ لاندې پوښتنه وګورئ۔ موږ پوهېږو
چې يو غير ناطق عدد ناطق ځواک ته پورته کول ناطق نتيجه ورکولے شي۔ د
بېلګې په توګه، $\sqrt{2}^2 = 2$۔ خو دا دومره څرګنده نۀ ده چې آيا يو
غير ناطق عدد \emph{غير ناطق} ځواک ته پورته کول هم ناطق نتيجه ورکولے
شي که نه۔ لاندې قضيه مثبت ځواب ورکوي:

\begin{thm}
داسې غير ناطق عددونه $a$ او $b$ شته چې $a^b$ ناطق وي۔
\end{thm}

\begin{proof}
$\sqrt{2}^{\sqrt{2}}$ وګورئ۔ که دا ناطق وي، نو خبره ثابته شوه:
$a = b = \sqrt{2}$ اخلو۔ که نه، نو دا غير ناطق دے۔ بيا لرو:
\[
(\sqrt{2}^{\sqrt{2}})^{\sqrt{2}} = \sqrt{2}^{\sqrt{2} \cdot
  \sqrt{2}} = \sqrt{2}^2 = 2,
\]
چې ناطق دے۔ نو په دې حالت کښې $a$ د
$\sqrt{2}^{\sqrt{2}}$ او $b$ د~$\sqrt 2$ په توګه اخلو۔
\end{proof}

آيا دا معتبر ثبوت دے؟ ډېری رياضيپوهان ئې معتبر ګڼي۔ خو بيا هم يوه
ناخوښه نکته شته: موږ د حقيقي عددونو د داسې يوې جوړې شتون ثابت کړے دے
چې ټاکلے خاصيت لري، خو دا نۀ شو وئيلے چې \emph{کومه} جوړه ده۔ همدا
پايله په داسې ډول هم ثابتولے شو چې د $a$ او $b$ جوړه \emph{په ثبوت کښې
ورکړل شوې} وي: $a = \sqrt{3}$ او $b = \log_3 4$ واخلئ۔ بيا:
\[
a^b = \sqrt{3}^{\log_3 4} = 3^{1/2 \cdot \log_3 4} = (3^{\log_3
  4})^{1/2} = 4^{1/2}= 2,
\]
ځکه چې $3^{\log_3 x} = x$۔

شهودي منطق د داسې استدلال د نيولو لپاره جوړ شوے چې د لومړي ثبوت په
شان ګامونه پکښې نۀ منل کېږي۔ د داسې $x$ شتون ثابتول چې~$!A(x)$ پوره
کوي، دا معنا لري چې بايد يو ځانګړے~$x$ او دا ثبوت ورکړئ چې $!A$ پوره
کوي، لکه په دويم ثبوت کښې۔ د $!A$ يا $!B$ ثابتول غواړي چې له دواړو
څخه يو ثابت کړے شئ۔

په رسمي ژبه، شهودي منطق د کلاسيک منطق د !!a{derivation} نظام له يوې
ټاکلې محدودونې څخه لاس ته راځي۔ له رياضيکي نظره دا يوازې رسمي استنتاجي
نظامونه دي؛ خو لکه چې وويل شول، موخه ئې د رياضيکي استدلال يو ډول نيول
دي۔ يو څوک دا هغه استدلال ګڼلے شي چې د رياضياتو د يو ځانګړي فلسفي
نظر، لکه د براور د شهوديت، له مخې موجه وي؛ بل څوک ئې د بېشپ د تعميريت
په شان لا «محسوس» او قانع کوونکے رياضيکي استدلال ګڼلے شي؛ او دا بحث
هم کېدے شي چې آيا رسمي تعريف غيررسمي انګېزه په سمه توګه نيسي که نه۔
خو هر فلسفي دريځ چې ولرو، شهودي منطق د يوه په رسمي توګه وړاندې شوي
منطق په توګه څېړلے شو؛ او ډېر رياضيکي منطقپوهان، د هر دليل له مخې،
په دې څېړنه کښې دلچسپي لري۔

\end{document}
